\section{El Proceso}

El desarrollo se articula en un pipeline unificado que automatiza desde la adquisición de datos hasta la validación de resultados:

\begin{itemize}
    \item \textbf{Selección e Integración:} Se unifican los datasets de \textit{LFW (control)} e \textit{Illinois DOC (riesgo)}. Para garantizar la imparcialidad, se han generado metadatos de etnia para LFW mediante el modelo \textbf{DeepFace}, permitiendo un \textit{muestreo estratificado (4-way balance)} que equilibra fotos de individuos blancos y negros de ambas fuentes, mitigando sesgos algorítmicos históricos.

    \item \textbf{Preprocesamiento y Limpieza (MTCNN):} Se utiliza la arquitectura \textbf{MTCNN (Multi-task Cascaded Convolutional Networks)} para detectar y extraer exclusivamente el óvalo facial de las imágenes originales, eliminando fondos, vestimenta o entornos que podrían confundir al clasificador. Las caras se reescalan a una resolución estándar de $224 \times 224$ píxeles y se normalizan con la media y desviación típica de ImageNet.

    \item \textbf{Minería de Datos (Entrenamiento):} Se ejecutan experimentos comparativos entre \textbf{ResNet50} y \textbf{VGG-16}. El entrenamiento se acelera mediante hardware especializado (\textit{MPS} en Apple Silicon o \textit{CUDA} en NVIDIA), implementando aumentación de datos mediante volteos horizontales aleatorios para mejorar la robustez.

    \item \textbf{Interpretación y Explicabilidad (Grad-CAM):} Se integra la técnica \textbf{Grad-CAM (Gradient-weighted Class Activation Mapping)} para generar mapas de calor sobre las fotos. Esto permite visualizar qué regiones del rostro (ojos, boca, mandíbula) están activando la predicción de "riesgo", aportando transparencia al "caja negra" del modelo.
\end{itemize}
